% ===== PRÉAMBULE CANONIQUE — à coller VERBATIM en tête de CHAQUE .tex =====
\documentclass[11pt,a4paper]{article}
\usepackage[utf8]{inputenc}\usepackage[T1]{fontenc}\usepackage{lmodern}
\usepackage[french]{babel}
\usepackage{amsmath,amssymb,mathtools}\usepackage{icomma}
\usepackage[version=4]{mhchem}          % \ce{...} : équations & formules
\usepackage{siunitx}\sisetup{locale=FR} % \qty{}{}, \si{}, \ang{}
\usepackage{chemfig}                    % \chemfig{...} : structures développées
\usepackage{xcolor}\usepackage{colortbl}
\usepackage{tikz}\usepackage{pgfplots}\pgfplotsset{compat=1.18}
\usepackage[most]{tcolorbox}\usepackage{enumitem}\usepackage{array,booktabs}
\usepackage[a4paper,margin=2.2cm]{geometry}
\usetikzlibrary{arrows.meta,calc,angles,quotes,positioning,patterns,backgrounds,shapes.geometric}
\definecolor{primary}{RGB}{222,49,99}\definecolor{primarydark}{RGB}{150,22,55}
\definecolor{defcol}{RGB}{0,114,178}\definecolor{methcol}{RGB}{52,92,148}
\definecolor{excol}{RGB}{0,135,90}\definecolor{attncol}{RGB}{200,40,55}
\tcbset{boxbase/.style={enhanced,breakable,boxrule=0.9pt,arc=2.5pt,
  left=3mm,right=3mm,top=2.3mm,bottom=2.3mm,fonttitle=\bfseries,coltitle=white}}
% --- boîtes TUTORIEL (noms EXACTS, ne pas renommer) ---
\newtcolorbox{cle}[1][]{boxbase,colback=defcol!6,colframe=defcol,
  title={\textsc{À retenir}\ifx\relax#1\relax\relax\else\textmd{ : \itshape #1}\fi}}
\newtcolorbox{methode}[1][]{boxbase,colback=methcol!7,colframe=methcol,sharp corners,
  title={\textsc{Méthode}\ifx\relax#1\relax\relax\else\textmd{ --- \itshape #1}\fi}}
\newtcolorbox{exemplebox}[1][]{boxbase,colback=excol!5,colframe=excol,
  title={\textsc{Exemple}\ifx\relax#1\relax\relax\else\textmd{ : \itshape #1}\fi}}
\newtcolorbox{piege}[1][]{boxbase,colback=attncol!6,colframe=attncol,
  title={\textsc{Piège}\ifx\relax#1\relax\relax\else\textmd{ : \itshape #1}\fi}}
\newtcolorbox{remarque}[1][]{boxbase,colback=black!4,colframe=black!45,
  title={\textsc{Remarque}\ifx\relax#1\relax\relax\else\textmd{ : \itshape #1}\fi}}
% --- boîtes EXERCICE / CORRIGÉ (noms EXACTS, auto-comptées) ---
\newtcolorbox[auto counter]{exo}[1][]{boxbase,colback=primary!4,colframe=primary,
  title={\textsc{Exercice}~\thetcbcounter\ifx\relax#1\relax\relax\else\textmd{ --- \itshape #1}\fi}}
\newtcolorbox[auto counter]{corrige}[1][]{boxbase,colback=excol!4,colframe=excol,
  title={\textsc{Corrigé}~\thetcbcounter\ifx\relax#1\relax\relax\else\textmd{ --- \itshape #1}\fi}}
\newcommand{\R}{\mathbb{R}}\newcommand{\N}{\mathbb{N}}\newcommand{\Z}{\mathbb{Z}}\newcommand{\ds}{\displaystyle}
% (un \usepackage{fancyhdr}+en-tête est possible ; il est ignoré côté web)
% =========================================================================
\begin{document}

\begin{corrige}[masse molaire par le volume molaire]
Aux CNTP, on passe par les moles.
\begin{enumerate}[label=\alph*)]
  \item $n=\dfrac{V}{V_{\mathrm{m}}}=\dfrac{1{,}12}{22{,}4}=\qty{0.05}{\mole}$.
  \item $M=\dfrac{m}{n}=\dfrac{2{,}2}{0{,}05}=\qty{44}{\gram\per\mole}$.
  \item $M=44$ correspond à \ce{CO2} ($12+2\times16$) ou \ce{N2O} ($2\times14+16$).
\end{enumerate}
\end{corrige}

\begin{corrige}[identification par masse volumique hors CNTP]
On part de $\rho=\dfrac{pM}{RT}$.
\begin{enumerate}[label=\alph*)]
  \item $M=\dfrac{\rho RT}{p}$.
  \item $M=\dfrac{1{,}30\times0{,}082\times300}{1{,}0}=\qty{31.98}{\gram\per\mole}\approx\qty{32}{\gram\per\mole}$.
  \item $M\approx32$ : c'est le dioxygène \ce{O2}.
\end{enumerate}
\end{corrige}

\begin{corrige}[un gaz plus léger que l'air]
\begin{enumerate}[label=\alph*)]
  \item $M=28{,}8\times d=28{,}8\times0{,}55\approx\qty{16}{\gram\per\mole}$.
  \item $M(\ce{CH4})=12+4=16$, $M(\ce{C2H6})=30$, $M(\ce{C3H8})=44$. C'est donc le méthane \ce{CH4}.
  \item $d=0{,}55<1$ : plus léger que l'air, le gaz \textbf{monte}.
\end{enumerate}
\end{corrige}

\begin{corrige}[le \ce{CO2} coule-t-il dans l'air ?]
\begin{enumerate}[label=\alph*)]
  \item $\rho(\ce{CO2})=\dfrac{M}{V_{\mathrm{m}}}=\dfrac{44}{22{,}4}=\qty{1.96}{\gram\per\litre}$.
  \item $d=\dfrac{\rho(\ce{CO2})}{\rho_{\text{air}}}=\dfrac{1{,}96}{1{,}29}\approx1{,}52$ (on retrouve
        $\dfrac{44}{28{,}8}=1{,}53$).
  \item $d>1$ : le \ce{CO2} est plus lourd que l'air, il s'accumule \textbf{au sol}. Dans un
        extincteur, il recouvre le foyer et chasse l'air (le dioxygène), étouffant la flamme.
\end{enumerate}
\end{corrige}

\end{document}
